\documentclass[a4paper,12pt]{article}

\usepackage[utf8]{inputenc}
\usepackage[T1]{fontenc}
\usepackage[italian]{babel}
\usepackage{enumitem}
\usepackage{geometry}
\usepackage{hyperref}
\usepackage{booktabs}
\usepackage{graphicx}
\usepackage{longtable}

\geometry{margin=2.5cm}

\title{Documentazione Tecnica di Progetto \\ \large Gestionale Locale per la Gestione del Magazzino}
\author{Gruppo 11}
\date{\today}

\begin{document}

\maketitle

\section*{Storico delle Versioni}
\begin{center}
\begin{tabular}{llll}
\toprule
\textbf{Versione} & \textbf{Data} & \textbf{Descrizione} & \textbf{Autore} \\
\midrule
v1.0 & 26/05/2026 & Strutturazione iniziale e analisi requisiti & Gruppo 11 \\
\bottomrule
\end{tabular}
\end{center}

\newpage
\tableofcontents
\newpage

\section{Introduzione}
\subsection{Scopo del Progetto}
Il software ha l'obiettivo di gestire le dinamiche logistiche interne di un'azienda, ottimizzando la gestione del magazzino, delle merci in entrata/uscita e della documentazione contabile associata.

\subsection{Obiettivi e Ambito (Scope)}
L'applicazione opera in un contesto puramente locale (offline). Questa scelta risponde a vincoli di sicurezza e alla necessità di operare in ambienti interni senza la necessità di una connessione a sistemi esterni. Il sistema gestisce l'intero ciclo di vita della merce: dal carico a magazzino alla generazione dei documenti di trasporto (DDT) e fatturazione.

\subsection{Glossario}
\begin{itemize}
    \item \textbf{DDT}: Documento di Trasporto.
    \item \textbf{RBAC}: Role-Based Access Control (Gestione permessi basata sui ruoli).
    \item \textbf{MVC}: Model-View-Controller (Pattern architetturale).
    \item \textbf{DAO}: Data Access Object (Oggetto per l'accesso ai dati).
    \item \textbf{Stock-out}: Condizione di esaurimento delle scorte.
\end{itemize}

\section{Analisi dei Requisiti}
\subsection{Attori del Sistema}
\begin{itemize}
    \item \textbf{Magazziniere}: Gestisce il flusso fisico delle merci, scansione codici a barre e generazione DDT.
    \item \textbf{Contabile}: Gestisce fatture, note di credito e listini prezzi.
    \item \textbf{Amministratore}: Supervisione globale, gestione utenti e reportistica avanzata.
\end{itemize}

\subsection{Requisiti Funzionali (RF)}
\begin{longtable}{p{0.15\textwidth} p{0.6\textwidth} p{0.15\textwidth}}
\caption{Requisiti Funzionali} \\
\toprule
\textbf{ID} & \textbf{Descrizione} & \textbf{Priorità} \\
\midrule
RF-01 & Login degli utenti tramite username e password (RBAC). & Alta \\
RF-02 & Visualizzazione e inserimento merci in entrata/uscita. & Alta \\
RF-03 & Generazione automatica DDT per merci in uscita. & Alta \\
RF-04 & Gestione giacenze e modifica manuale delle quantità. & Alta \\
RF-05 & Notifica visiva per articoli sotto la soglia minima. & Media \\
RF-06 & Gestione e consultazione resi e storico ordini. & Media \\
RF-07 & Generazione fatture e note di credito. & Alta \\
RF-08 & Modifica listini prezzi (acquisto/vendita). & Media \\
RF-09 & Scansione codici a barre per carico/scarico rapido. & Media \\
\bottomrule
\end{longtable}

\subsection{Requisiti Non Funzionali (RNF)}
\begin{longtable}{p{0.15\textwidth} p{0.25\textwidth} p{0.45\textwidth}}
\caption{Requisiti Non Funzionali} \\
\toprule
\textbf{ID} & \textbf{Categoria} & \textbf{Descrizione} \\
\midrule
RNF-01 & Prestazioni & Tempo di risposta scansione < 0.5s. \\
RNF-02 & Prestazioni & Generazione report annuale < 3s. \\
RNF-03 & Portabilità & Funzionamento offline su sistemi desktop moderni. \\
RNF-04 & Sicurezza & Dati archiviati localmente senza esposizione esterna. \\
RNF-05 & Archiviazione & Database embedded SQLite (singolo file). \\
RNF-06 & Usabilità & Interfaccia GUI leggera e reattiva. \\
\bottomrule
\end{longtable}

\subsection{Analisi dei Casi d'Uso}
Il sistema è documentato tramite diagrammi dei casi d'uso (rif. \texttt{Use\_case\_diagram.puml}) che mappano le interazioni tra i tre attori principali e le macro-funzionalità del sistema (Login, Gestione Ordini, Gestione Giacenze).

\section{Progettazione e Architettura}
\subsection{Scelte Architetturali}
Il software adotta il pattern \textbf{MVC (Model-View-Controller)}. Questa scelta garantisce la separazione tra logica di business e interfaccia utente. Il database scelto è \textbf{SQLite} per garantire portabilità assoluta (zero configurazione) e facilità di backup (copia di un singolo file).

\subsection{Modello del Dominio}
Il dominio è centrato sulle entità \texttt{Articolo}, \texttt{Movimento}, \texttt{Ordine} e \texttt{Utente}. (Rif. \texttt{Domain\_model.puml}).

\subsection{Design Dettagliato (Diagramma delle Classi)}
La struttura delle classi segue il pattern DAO per l'accesso ai dati, garantendo che il Controller non dipenda direttamente dalle query SQL. (Rif. \texttt{Class\_diagram.puml}).

\section{Gestione del Progetto}
\subsection{Processo di Sviluppo}
Viene utilizzato un approccio Agile ispirato a \textbf{Scrum}, con un Backlog prioritizzato che include la configurazione dell'ambiente, lo sviluppo dei moduli CRUD e l'integrazione del lettore barcode.

\subsection{User Stories}
\textbf{US-01}: «Come magazziniere, voglio scansionare il codice a barre di un articolo, per aggiornare la giacenza istantaneamente».
\begin{itemize}
    \item \textit{Criterio 1}: Riconoscimento standard EAN-13.
    \item \textit{Criterio 2}: Aggiornamento DB automatico post-scansione.
\end{itemize}

\section{Strategia di Verifica e Validazione}
\subsection{Piano di Testing}
Il piano prevede Unit Testing sistematico sulla Business Logic tramite JUnit. Particolare attenzione è rivolta alla gestione delle transazioni.

\subsection{Casi di Test (Test Cases)}
\textbf{TC-01}: Tentativo di scarico merce con quantità superiore alla giacenza.
\begin{itemize}
    \item \textit{Input}: Scarico 6 unità (Giacenza: 5).
    \item \textit{Risultato Atteso}: \texttt{GiacenzaInsufficienteException} e Rollback transazione.
\end{itemize}

\end{document}
